% ============================================================
% s03-curve-b.tex —— §1.5 概形曲线、§1.6 untilt 作为 Cartier 除子
% 底本：Laurent Fargues, Geometrization of the local Langlands
%       correspondence: an overview（2016-02-02），原文 p.10–p.16
% 本文件承接 s02-curve-a（§1.1–§1.4），仍属原文第 1 节，故不写
% \section；原文 §2 起见 s04-bung.tex。
% ============================================================


% —— 单章草稿自检用：仅当此前没有 \section 时预置编号，
%    使定理环境从原文 1.24 起排。全书编译时 section=1，此块自动跳过。 ——
\ifnum\value{section}=0\relax
  \setcounter{section}{1}%
  \setcounter{theorem}{23}%
  \setcounter{equation}{5}%
\fi

\subsection{概形曲线}% 原文 1.5. The schematical curve
\srcpage{10}

\subsubsection*{1.5.1.\ 定义}
\srcpage{11}

设 $S=\Spa(R,R^+)$ 是 $\FF_q$ 上的 affinoid 完满空间。对每个 $d\in\ZZ$，$Y_{R,R^+}$ 上有一个线丛 $\oo(d)$，其几何实现为
\[
Y\times_{\varphi^{\ZZ}}\AA^1\longrightarrow Y/\varphi^{\ZZ},
\]
其中 $\varphi$ 通过 $\times\pi^{-d}$ 作用在 $\AA^1$ 上。现在我们借助下述定义宣布 $\oo(1)$ 是丰沛的。

\begin{definition}\label{def:1.24}
(1) 令
\[
P_R=\bigoplus_{d\ge 0}H^0(X_{R,R^+},\oo(d))=\bigoplus_{d\ge 0}\oo(Y_{R,R^+})^{\varphi=\pi^d}.
\]
\yizhu{$P_R$ 以 $d$ 为分次；$\oo(Y)^{\varphi=\pi^d}$ 表示满足 $\varphi(x)=\pi^d x$ 的元素，它正是 $H^0$ 的显式描述。分次环既已给定，便可取 Proj。}{A：原文 §1.5.1 p.11}

(2) 定义 $X_R=\Proj(P_R)$。当有歧义时，我们用记号 $X_R^{sch}$ 与 $X^{ad}_{R,R+}$ 指明所用的究竟是概形曲线还是 adic 曲线。\yizhu{“概形曲线” 即原文的 schematical curve：由分次环 $P_R$ 的 Proj 得到的概形；与之相对的 “adic 曲线” $X^{ad}$ 是 adic 空间。两者由态射 $X^{ad}\to X^{sch}$ 相联系。}{A：原文 §1.5.1 p.11}
\end{definition}

$R$ 是完满域 $F$ 的情形正是 [10] 的研究对象。\emph{此时 $X_F$ 是 Dedekind 概形。}一般地，它不满足任何有限性性质。

恒有环空间的态射
\[
X^{ad}_{R,R^+}\longrightarrow X^{sch}_R .
\]

\subsubsection*{1.5.2.\ 一个等特征下的计算}

也许值得更详细地解释特征 $p$ 的情形。于是设 $E=\FF_q((\pi))$。此时有 $\FF_q$-向量空间的同构
\[
\begin{aligned}
(R^{\circ\circ})^d&\xrightarrow{\ \sim\ }\oo(Y_{R,R^+})^{\varphi=\pi^d},\\
(x_0,\dots,x_{d-1})&\longmapsto\sum_{i=0}^{d-1}\sum_{n\in\ZZ}x_i^{q^{-n}}\pi^{nd+i}.
\end{aligned}
\]

考虑 $\FF_q$ 上的形式群 $\mathscr{G}_d=\widehat{\mathbb{G}}_a^d$，它带有 $\oo_E=\FF_q\llbracket\pi\rrbracket$ 的作用，其中 $\pi$ 的作用为
\[
(X_0,\dots,X_{d-1})\longmapsto(X_{d-1}^q,X_0,\dots,X_{d-2}).
\]
那么前述同构是 $E$-向量空间的同构
\[
\mathscr{G}_d(R^\circ)\xrightarrow{\ \sim\ }\oo(Y_{R,R^+})^{\varphi=\pi^d}.
\]

现在考虑幂级数
\[
f(T)=\sum_{n\ge 0}\frac{T^{q^n}}{\pi^n}\in E\llbracket T\rrbracket .
\]
可以验证
\[
f^{-1}(\pi f)\in\oo_E\llbracket T\rrbracket ,
\]
其中 $f^{-1}$ 是复合意义下的逆；它还满足
\[
f^{-1}(\pi f)\equiv\pi T+\alpha T^q\ \ \mathrm{mod}\ \ T^{q^2}\oo_E\llbracket T\rrbracket ,
\]
其中 $\alpha\in 1+\pi\oo_E$。把 $\mathscr{G}_1$ 记作 $\mathscr{G}$。于是 $\mathscr{G}$ 在 $\oo_E$ 上有提升 $\widetilde{\mathscr{G}}$，其底形式群是 $\widehat{\mathbb{G}}_a$，且 $\pi$ 的作用由
\[
[\pi]_{\widetilde{\mathscr{G}}}=f^{-1}(\pi f)
\]
给出，也就是说
\[
f=\log_{\widetilde{\mathscr{G}}}.
\]

取伪一致化元 $\varepsilon\in R^{\circ\circ}\cap R^\times$。我们要找
\[
t_\varepsilon:=\sum_{n\in\ZZ}\varepsilon^{q^{-n}}\pi^n\in\oo(Y_{R,R^+})^{\varphi=\pi}
\]
的 Weierstrass 乘积展开。

\leavevmode\srcpage{12}定义
\[
\eta=\lim_{k\to+\infty}[\pi^k]_{\widetilde{\mathscr{G}}}(\varepsilon^{q^{-k}})\in R^{\circ\circ}\llbracket\pi\rrbracket ,
\]
其中极限取 $\mathbf{A}=R^\circ\llbracket\pi\rrbracket$ 的 $(\varpi_R,\pi)$-进拓扑。再令
\[
u_\varepsilon=\frac{\eta}{\varphi^{-1}(\eta)}=\frac{[\pi]_{\widetilde{\mathscr{G}}}(\varphi^{-1}(\eta))}{\varphi^{-1}(\eta)}\in\mathbf{A}.
\]

可以验证 $u_\varepsilon$ 是次数 1 的本原元，而且还有
\[
u_\varepsilon=\varepsilon^{\frac{q-1}{q}}+\pi(1+\beta),
\]
其中 $\beta\in R^{\circ\circ}\llbracket\pi\rrbracket$。于是可以构造收敛的 Weierstrass 乘积
\[
\Pi^+(u_\varepsilon)=\prod_{k\ge 0}\frac{\varphi^k(u_\varepsilon)}{\pi}\in\oo(Y_{R,R^+}).
\]

现在令
\[
\Pi^-(u_\varepsilon)=\varphi^{-1}(\eta),
\]
可以把它想象成 “$\prod_{k<0}\varphi^k(u_\varepsilon)$”，尽管这个乘积并不收敛。即使这个乘积没有意义，我们仍要指出有部分分解
\[
\eta=\varphi^{i-1}(\eta)\cdot\prod_{i\le k<0}\varphi^k(u_\varepsilon)
\]
对任意 $i<0$ 成立。我们有
\[
\Pi(u_\varepsilon)=\Pi^+(u_\varepsilon)\Pi^-(u_\varepsilon)\in\oo(Y_{R,R^+})^{\varphi=\pi}
\]
以及
\[
\Pi(u_\varepsilon)=t_\varepsilon .
\]

最后这个等式来自刚性解析函数 $\log_{\widetilde{\mathscr{G}}}$ 作为 $\mathbb{B}^1_E$ 上全纯函数的 Weierstrass 乘积展开：
\[
\log_{\widetilde{\mathscr{G}}}=\lim_{k\to+\infty}\frac{[\pi^k]_{\widetilde{\mathscr{G}}}}{\pi^k}.
\]

考虑
\[
R_\varepsilon=\mathbf{A}\Big[\frac{1}{\pi}\Big]\Big/u_\varepsilon ,
\]
即与次数 1 本原元 $u_\varepsilon$ 相联系的 $R$ 的 “untilt”。由于我们在等特征情形，untilt 无非是给出一个连续态射 $E\to R$；也就是说，对 $u_\varepsilon$ 作 Weierstrass 分解之后其实有 $R_\varepsilon=R$，只是我们不想固定这样的等同。记
\[
\theta_\varepsilon:\mathbf{A}\Big[\frac{1}{\pi}\Big]\longrightarrow R_\varepsilon
\]
为投影。那么 $\big(\theta_\varepsilon(\varphi^{-n}(\eta))\big)_{n\ge0}$ 是 Tate 模 $T_\pi(\widetilde{\mathscr{G}}\otimes_{\oo_E}R_\varepsilon^\circ)$ 的一组基。

现在设 $I\subset\,]0,1[$ 是端点落在 $p^{\mathbb{Q}}$ 中的区间，并考虑 $D^*_{(R,R+),I}$，即 $D^*_{R,R+}$ 中经由函数 $\delta:|Y_{R,R+}|\to\,]0,1[$ 对应的圆环，其中 $\delta$ 在固定某个 $\varpi_R$ 之后定义。我们固定 $\|\varpi_R\|=1/p$。

\begin{definition}\label{def:1.25}
若 $\xi=\sum_{n\ge0}a_n\pi^n$ 是次数 1 的本原元，则定义 $\|\xi\|=\|a_0\|$。
\end{definition}

我们指出：若 $u\in R^\circ\llbracket\pi\rrbracket^\times$，则 $\|u\xi\|=\|\xi\|$。因此，若 $\xi=u(\pi-a)$ 是 $\xi$ 的 Weierstrass 分解，则 $\|\xi\|=\|a\|$。

让我们回到前面的情形。我们有 $\|u_\varepsilon\|=\|\varepsilon\|^{\frac{q-1}{q}}$。于是可以验证
\begin{itemize}[leftmargin=2em]
\item 若 $A=\{k\in\ZZ\mid \|u_{\varepsilon^{q^k}}\|\in I\}$，则 $t_\varepsilon=\prod_{k\in A}\varphi^k(u_\varepsilon)\times u$，其中 $u\in\oo(D^*_I)^\times$；
\item 对任意 $n\ge2$，$\oo(D^*)u_\varepsilon+\oo(D^*)u_{\varepsilon^n}=\oo(D^*)$。
\end{itemize}

\leavevmode\srcpage{13}由此推出：若 $(n,p)=1$，则
\[
\oo(D^*_I)t_\varepsilon+\oo(D^*_I)t_{\varepsilon^n}=\oo(D^*_I).
\]

我们证明了下面这一点。

\begin{proposition}\label{prop:1.26}
存在 $\varepsilon_1,\varepsilon_2\in R^{\circ\circ}\cap R^\times$，使得
\[
\oo(D^*_I)t_{\varepsilon_1}+\oo(D^*_I)t_{\varepsilon_2}=\oo(D^*_I).
\]
\end{proposition}

于是有覆盖
\[
\Spec(\oo(D^*_I))=\bigcup_{\varepsilon\in R^{\circ\circ}\cap R^\times}D(t_\varepsilon).
\]

对这样的 $\varepsilon$，有态射
\[
\Spec\Big(\oo(D^*_I)\Big[\frac{1}{t_\varepsilon}\Big]\Big)\longrightarrow D^+(t_\varepsilon)\subset\Proj(P).
\]

当 $\varepsilon$ 与 $I$ 变动时，这给出 ind-概形\yizhu{ind-scheme 指归纳极限仍是概形的一类对象，此处仍写作 “ind-概形”，与工程内 adic、pro-étale 等前缀词的处理一致。}{I：本工程术语惯例；术语表暂无对应条目}的 $\varphi$-不变态射
\[
\varinjlim_{I}\Spec(\oo(D^*_I))\longrightarrow\Proj(P)=X^{sch}.
\]

经由环空间的态射
\[
D^*_I\longrightarrow\Spec(\oo(D^*_I))
\]
（在拓扑空间层面上是赋值的支撑），它诱导出态射
\[
X^{ad}\longrightarrow X^{sch}.
\]

\subsubsection*{1.5.3.\ 不等特征下的类似计算}
% 原文标题印作 “The analog computation in equal characteristic”，实为不等特征，见下方译注。

前一节几乎所有结果都可推广到 $[E:\QQ_p]<+\infty$ 的情形。\yizhuw{%
原文标题印作“The analog computation in equal characteristic”，但本节通篇假设
$[E:\QQ_p]<+\infty$（不等特征），真正的等特征计算在 §1.5.2。
标题显系笔误，译文按内容改题；原文误标见 \texttt{工程文件/work/疑点清单.tsv}。}
这里的 “几乎” 是因为：当 $d\ge2$ 时，$(R,R^+)\mapsto\oo(Y_{R,R^+})^{\varphi=\pi^d}$ 不能由形式群（更确切地说，其万有覆盖）表示。但对 $d=1$，若取 $\mathscr{G}$ 为 $\oo_E$ 上以
\[
\log_{\mathscr{G}}=\sum_{n\ge0}\frac{T^{q^n}}{\pi^n}
\]
为对数的 Lubin–Tate 群律，其归约为 $\mathscr{G}$，则
\[
\begin{aligned}
\mathscr{G}(R^\circ)&\xrightarrow{\ \sim\ }\oo(Y_{R,R^+})^{\varphi=\pi},\\
\varepsilon&\longmapsto t_\varepsilon:=\sum_{n\in\ZZ}[\varepsilon^{q^{-n}}]\pi^n .
\end{aligned}
\]

此外，对 $\varepsilon\in R^{\circ\circ}\cap R^\times$，$t_\varepsilon$ 有 Weierstrass 乘积展开
\[
t_\varepsilon=\Pi^-(u_\varepsilon)\cdot\Pi^+(u_\varepsilon),
\]
其中
\[
u_\varepsilon=\frac{\eta}{\varphi^{-1}(\eta)}
\]
并且
\[
\eta=\lim_{k\to+\infty}[\pi^k]_{\mathscr{G}}\big([\varepsilon^{q^{-k}}]\big).
\]

所有这些使我们能构造环空间的态射
\[
X^{ad}_{R,R+}\longrightarrow X^{sch}_R .
\]

\srcpage{14}

\subsubsection*{1.5.4.\ 概形曲线的内蕴定义与 GAGA}

我们把 $X^{sch}$ 定义为分次代数 $P$ 的 Proj。于是它典范地带有一个线丛 $\oo(1)=\widetilde{P}[1]$。但事实上，没有理由固定这样一个线丛：\emph{曲线本身存在，其上并没有固定的线丛 $\oo(1)$}。例如，我们所取的 $\oo(1)$ 的定义依赖于一致化元 $\pi$ 的选择，换一种选择可能得到不同构的 $\oo(1)$。下一节还会清楚地看到，$X^{sch}$ 上有大量与 $R$ 的 untilt 相联系的丰沛线丛可选；前一节所取的那个，对应于满足 $T_\pi(\widetilde{\mathscr{G}}\otimes_{\oo_E}R^{\sharp,\circ})=\oo_E$ 的 untilt $R^\sharp$（这一条件依赖于 $\pi$ 的选择，因为按定义 $\widetilde{\mathscr{G}}$ 就依赖于该选择）。

尽管如此，我们有下述命题，其证明并不很难（底域的情形见 [10] 第 6.7 节）。

\begin{proposition}\label{prop:1.27}
ind-概形的 $\varphi$-不变态射
\[
\varinjlim_{I}\Spec(\oo(Y_I))\longrightarrow X^{sch}
\]
使 $X^{sch}$ 成为 ind-概形 $\varinjlim_I\Spec(\oo(Y_I))$ 在 $\varphi^{\ZZ}$ 作用下的范畴商。
\end{proposition}

\begin{corollary}\label{cor:1.28}
为使概形曲线 $X^{sch}$ 成为内蕴的，我们必须改变它的定义：$X^{sch}$ 是 $\varinjlim_I\Spec(\oo(Y_I))$ 在 $\varphi^{\ZZ}$ 作用下的范畴商。
\end{corollary}

当然，要证明这样的范畴商存在，最好的办法就是像我们这样取定义 1.24。

后面我们会用到下述定理。

\begin{theorem}[{[16]}]\label{thm:1.29}
与环空间态射 $X^{ad}_{R,R+}\to X^{sch}_R$ 相联系的 GAGA 函子是等价
\[
\Bun_{X_R}\xrightarrow{\ \sim\ }\Bun_{X^{ad}_{R,R+}}.
\]
\end{theorem}

\subsection{untilt 作为概形曲线上的 Cartier 除子}% 原文 1.6. Untilts as Cartier divisors on the schematical curve

设 $S\in\Perf_{\FF_q}$。如我们所见，$S$ 的任一 untilt $S^\sharp$ 给出 Cartier 除子
\[
S^\sharp\hookrightarrow Y_S .
\]

若 $S=\Spa(R,R^+)$ 是 affinoid 完满空间，则这个 Cartier 除子由可逆理想 $\oo_{Y_S}.\xi\subset\oo_{Y_S}$ 定义，其中 $\xi$ 是与 $S^\sharp$ 相联系的本原元。

这个 Cartier 除子 $D$ 给出 $\varphi$-不变的 Cartier 除子
\[
\bigoplus_{k\in\ZZ}\varphi^{k*}D,
\]
也就是说 $X^{ad}_S$ 上的一个 Cartier 除子。若 $S=\Spa(R,R^+)$，按 GAGA，它必须来自 $X^{sch}$ 上的某个 Cartier 除子。下面就是它的描述。

\subsubsection*{1.6.1.\ 等特征情形}

设 $\xi\in R^\circ\llbracket\pi\rrbracket$ 是次数 1 的本原元。与之相联系的 $Y=D^*_{R,R+}$ 上的 Cartier 除子，由线丛 $\xi^{-1}\oo_Y$ 连同其截面 $1\in\oo(Y)$ 给出。

设 $\cL_\xi$ 是 $X^{ad}$ 上相应的线丛。可以算出，对任意 $d\ge0$
\begin{equation}\label{eq:1.6}
\Gamma(X^{ad},\cL_\xi(d))\subset\oo(Y)\Big[\frac{1}{\xi}\Big]^{\varphi=\pi^d},
\end{equation}
而 $\cL_\xi^{ad}$ 的、定义 $X^{ad}$ 上该 Cartier 除子的截面，由 $1\in\oo(Y)\big[\frac{1}{\xi}\big]^{\varphi=\mathrm{Id}}$ 给出。把 $\xi$ 乘以 $R^\circ\llbracket\pi\rrbracket^\times$ 中的元素，我们可以假设 $\xi=\pi-a$，其中 $a\in R^{\circ\circ}\cap R^\times$。我们构造\srcpageinl{15}收敛的 Weierstrass 乘积
\[
\Pi^+(\xi):=\prod_{n\ge0}\frac{\varphi^n(\xi)}{\pi}=\prod_{n\ge0}\Big(1-\frac{a^{q^n}}{\pi}\Big).
\]
它满足函数方程
\[
\varphi\big(\Pi^+(\xi)\big)=\frac{\xi}{\pi}\cdot\Pi^+(\xi).\yizhu{%
按本页对 $\Pi^+$ 的定义直接计算，只能得到 $\varphi(\Pi^+(\xi))=\frac{\pi}{\xi}\Pi^+(\xi)$，
与印文互为倒数。二者必有一误；译文照排，未擅改。% 
}{A：原文 p.15 定义与同页函数方程；待核 [10]}
\]
于是可以验证\yizhuw{上标 $\xi\varphi=\pi^{d-1}$ 是原文的简写：先在 $\xi$ 处局部化（即 $\oo(Y)[1/\xi]$），再取其中满足 $\varphi(x)=\pi^{d-1}x$ 的元素。}：经由 (6)，
\[
\times\Pi^+(\xi)\colon\Gamma(X^{ad},\cL_\xi^{ad}(d))\xrightarrow{\ \sim\ }\oo(Y)^{\xi\varphi=\pi^{d-1}}.
\]

现在定义 $P=\bigoplus_{d\ge0}\oo(Y)^{\varphi=\pi^d}$ 上的分次模
\[
N_\xi=\bigoplus_{d\ge0}\oo(Y)^{\xi\varphi=\pi^{d-1}}.
\]
于是可以验证下面这一点。

\begin{proposition}\label{prop:1.30}
拟凝聚层 $\cL_\xi=\widetilde{N_\xi}$ 是 $X^{sch}=\Proj(P)$ 上的线丛。此外，赋予截面 $\Pi^+(\xi)\in\Gamma(X^{sch},\cL_\xi)$，它定义了 $X^{sch}$ 上的一个 Cartier 除子，该除子经 GAGA 对应于 $X^{ad}$ 上由 $\xi$ 定义的 Cartier 除子。
\end{proposition}

\subsubsection*{1.6.2.\ 不等特征情形}

设 $\xi\in\mathbf{A}=W_{\oo_E}(R^\circ)$ 是次数 1 的本原元。若 $\widetilde{R}=R^\circ/R^{\circ\circ}$，则按定义，$\xi$ 在 $W_{\oo_E}(\widetilde{R})$ 中的归约 $\widetilde{\xi}$ 可以写成 $\pi$ 乘一个单位。由此推出：把 $\xi$ 乘以 $\mathbf{A}$ 中的单位之后，可以假设 $\widetilde{\xi}=\pi$，也就是说
\[
\xi\in\pi+W_{\oo_E}(R^{\circ\circ}).
\]

我们于是可以构造收敛的 Weierstrass 乘积
\[
\Pi^+(\xi):=\prod_{n\ge0}\frac{\varphi^n(\xi)}{\pi}\in\oo(Y).
\]
我们于是有命题 1.30 的类似物。

\begin{proposition}\label{prop:1.31}
拟凝聚层 $\cL_\xi=\Big(\bigoplus_{d\ge0}\oo(Y)^{\xi\varphi=\pi^{d-1}}\Big)^{\sim}$ 是 $X^{sch}_R$ 上的线丛。它带有截面 $\Pi^+(\xi)\in\Gamma(X^{sch},\cL_\xi)$。二元组 $(\cL_\xi,\Pi^+(\xi))$ 定义 $X^{sch}_R$ 上的 Cartier 除子，它经 GAGA 对应于 $X^{ad}_{R,R+}$ 上与 $\xi$ 所定义的 $R$ 的 untilt 相联系的 Cartier 除子。
\end{proposition}

\subsubsection*{1.6.3.\ 沿 Cartier 除子的形式完备化}

设 $R$ 是完满 $\FF_q$-代数。设 $R^\sharp$ 是 $R$ 的 untilt，$\xi\in\mathbf{A}$ 是与之相联系的次数 1 本原元。

\begin{definition}\label{def:1.32}
记 $\mathrm{B}^{+}_{dR}(R^\sharp)$ 为 $\mathbf{A}\big[\frac{1}{\pi}\big]$ 的 $\xi$-进完备化。记 $\mathrm{B}_{dR}(R^\sharp)=\mathrm{B}^{+}_{dR}(R^\sharp)\big[\frac{1}{\xi}\big]$\yizhu{$\mathrm{B}^{+}_{dR}$ 即 Fontaine 的 de Rham 周期环的正部分。此处它被实现为 $\mathbf{A}[1/\pi]$ 沿 $\xi$ 的形式完备化，这正是原文引出的方式。}{A：原文 §1.6.3 p.15}。
\end{definition}

于是有下面这一点。

\begin{proposition}\label{prop:1.33}
$X^{sch}_R$ 沿 $R^\sharp$ 所定义 Cartier 除子的形式完备化是 $\Spf\big(\mathrm{B}^{+}_{dR}(R^\sharp)\big)$。
\end{proposition}

照例，等特征情形要简单得多。\yizhuw{原文作 “A usual”，当为 “As usual”（照例、如常）之误；译文按后者理解。}此时 $R^\sharp$ 由一个态射 $\FF_q((\pi))\to R$ 给出，$\pi\mapsto a$，并且可以取 $\xi=\pi-a$。那么 $\mathrm{B}^{+}_{dR}(R^\sharp)$ 就是 $R^\circ\llbracket\pi\rrbracket\big[\frac{1}{\pi}\big]$ 的 $\xi$-进完备化。于是有典范截面
\[
\begin{tikzcd}
\mathrm{B}^{+}_{dR}(R^\sharp)\arrow[r,two heads,"\theta"] & R^\sharp\\
& R\arrow[ul]\arrow[u,equal]
\end{tikzcd}
\]

\leavevmode\srcpage{16}这给出典范同构
\[
\begin{aligned}
R\llbracket T\rrbracket&\xrightarrow{\ \sim\ }\mathrm{B}^{+}_{dR}(R^\sharp),\\
T&\longmapsto\xi .
\end{aligned}
\]
但这里要小心一点：这是 $E$-代数的同构，其中 $R\llbracket T\rrbracket$ 带有由
\[
\begin{aligned}
\FF_q((\pi))&\longrightarrow R\llbracket T\rrbracket,\\
\pi&\longmapsto T+a
\end{aligned}
\]
给出的 $E$-代数结构。

% —— 本块结束。原文 §2 The stack Bun_G 起由 s04-bung.tex 承接。 ——
